% =====================================================================
%  Informe + Anexos — Actividad 1: LCS
%  Documento autocontenido: todo el código y los datos están incrustados.
%  Repositorio: https://github.com/carpediemservus/lcs-actividad1
%  Compilar con:
%      pdflatex informe.tex && pdflatex informe.tex
% =====================================================================
\documentclass[11pt,a4paper]{article}

% ---------- Codificación e idioma ----------
\usepackage[utf8]{inputenc}
\usepackage[T1]{fontenc}
\usepackage{lmodern}
\usepackage[spanish,es-nodecimaldot,es-tabla]{babel}

% ---------- Geometría y maquetación ----------
\usepackage[margin=2.2cm]{geometry}
\usepackage{amsmath}
\usepackage{booktabs}
\usepackage{array}
\usepackage{graphicx}
\usepackage{tikz}
\usepackage{pgfplots}
\pgfplotsset{compat=1.17}
\usepackage{caption}
\captionsetup{labelfont=bf,skip=4pt}

% ---------- Hipervínculos ----------
\usepackage{xcolor}
\usepackage[hidelinks]{hyperref}

% ---------- Listados de código ----------
\usepackage{listings}
\definecolor{codegreen}{rgb}{0,0.5,0}
\definecolor{codegray}{rgb}{0.5,0.5,0.5}
\definecolor{codeblue}{rgb}{0,0,0.6}
\definecolor{codered}{rgb}{0.6,0,0}
\definecolor{backcolour}{rgb}{0.97,0.97,0.97}

\renewcommand{\lstlistingname}{Listado}
\renewcommand{\lstlistlistingname}{Índice de listados}
\lstdefinestyle{py}{
  language=Python,
  backgroundcolor=\color{backcolour},
  commentstyle=\color{codegreen}\itshape,
  keywordstyle=\color{codeblue}\bfseries,
  stringstyle=\color{codered},
  numberstyle=\tiny\color{codegray},
  basicstyle=\ttfamily\normalsize,
  breaklines=true, breakatwhitespace=true,
  numbers=left, numbersep=6pt,
  showstringspaces=false, tabsize=4,
  frame=single, framerule=0.3pt,
  rulecolor=\color{codegray!50},
  xleftmargin=1em, xrightmargin=1em,
}
\lstdefinestyle{mk}{
  language=make,
  backgroundcolor=\color{backcolour},
  basicstyle=\ttfamily\normalsize, breaklines=true,
  numbers=left, numbersep=6pt, frame=single, framerule=0.3pt,
  rulecolor=\color{codegray!50},
  xleftmargin=1em, xrightmargin=1em,
}
\lstdefinestyle{txt}{
  backgroundcolor=\color{backcolour},
  basicstyle=\ttfamily\normalsize, breaklines=true,
  numbers=none, frame=single, framerule=0.3pt,
  rulecolor=\color{codegray!50},
  xleftmargin=1em, xrightmargin=1em,
}

% ---------- Ajustes tipográficos ----------
\setlength{\parskip}{2pt}
\setlength{\parindent}{0pt}

% ---------- Título ----------
\title{\vspace{-1.5cm}\textbf{Diseño, análisis y comparación de algoritmos\\para la Subsecuencia Común Más Larga}}
% >>> Completar con los datos del estudiante <<<
\author{Nombre y apellidos del estudiante \\ Actividad 1 · Análisis de Algoritmos}
\date{Fecha de entrega: \today}

\begin{document}
\maketitle
\vspace{-0.6cm}

% =====================================================================
%  RESUMEN
% =====================================================================
\begin{abstract}
\normalsize\noindent
Se presentan y comparan tres aproximaciones al problema de la Subsecuencia
Común Más Larga (LCS): una versión recursiva directa, la estrategia de
divide y vencerás de Hirschberg y una solución mediante programación
dinámica tabular. Se analiza su complejidad temporal y espacial, y se
contrastan experimentalmente sobre cinco parejas de cadenas de ADN de
longitud creciente hasta 1000 símbolos. Los resultados confirman el
comportamiento exponencial de la versión recursiva, el uso lineal de
memoria de Hirschberg y el equilibrio entre simplicidad y consumo de
memoria de la versión tabular. El código fuente completo, el protocolo
experimental y los resultados brutos se incluyen en los anexos de este
mismo documento y en el repositorio público
\url{https://github.com/carpediemservus/lcs-actividad1}.

\medskip
\noindent\textbf{Palabras clave:} LCS, divide y vencerás, Hirschberg,
programación dinámica, análisis de complejidad.
\end{abstract}

% =====================================================================
%  1. INTRODUCCIÓN
% =====================================================================
\section{Introducción}

Dadas dos cadenas $X$ e $Y$ de longitudes $m$ y $n$, la
\emph{Subsecuencia Común Más Larga} (LCS) es una subsecuencia de
longitud máxima que aparece en ambas cadenas conservando el orden
relativo de los símbolos, aunque no necesariamente de forma contigua.
El problema tiene aplicaciones directas en bioinformática (alineamiento
de ADN y proteínas), en la comparación de versiones de ficheros
(\texttt{diff}), en correctores ortográficos y en detección de plagio.

El objetivo de esta actividad es poner en práctica tres estrategias
algorítmicas clásicas sobre el mismo problema, medir empíricamente sus
diferencias de tiempo y memoria, y justificar las elecciones de diseño
en función de sus respectivas complejidades. Las tres versiones
implementadas devuelven la propia subsecuencia, no solo su longitud:

\begin{itemize}
    \item \textbf{Versión A} — recursión directa, sin memorización.
    \item \textbf{Versión B} — divide y vencerás según el algoritmo de
          Hirschberg.
    \item \textbf{Versión C} — programación dinámica tabular.
\end{itemize}

El resto del documento se organiza como sigue. La
Sección~\ref{sec:metodologia} describe el diseño de cada algoritmo, su
análisis de complejidad y el protocolo experimental. La
Sección~\ref{sec:resultados} recoge las mediciones de tiempo de CPU y
memoria máxima. La Sección~\ref{sec:discusion} interpreta los resultados
y la Sección~\ref{sec:conclusiones} cierra con las conclusiones. Los
anexos contienen el código fuente completo, los datos empleados y las
instrucciones de reproducción local.

% =====================================================================
%  2. METODOLOGÍA
% =====================================================================
\section{Metodología}\label{sec:metodologia}

\subsection{Definición recurrente del problema}

Sea $L(i,j)$ la longitud de la LCS de los prefijos $X[0..i)$ e $Y[0..j)$.
Se cumple:
\[
L(i,j)=
\begin{cases}
0, & i=0\ \text{o}\ j=0,\\[2pt]
L(i-1,j-1)+1, & X[i-1]=Y[j-1],\\[2pt]
\max\{L(i-1,j),\,L(i,j-1)\}, & \text{en otro caso}.
\end{cases}
\]

\subsection{Versión A — recursión directa}

Se traduce literalmente la recurrencia anterior a una función recursiva
que devuelve la subsecuencia de cada subproblema: si los símbolos
coinciden se añade el símbolo a la solución de $(i-1,j-1)$, y si no se
elige la más larga de las dos ramas. No se almacena ningún resultado
intermedio, por lo que cada subproblema se recalcula un número
exponencial de veces. Su estado adicional es la pila de recursión y las
cadenas parciales de cada nivel.

\subsection{Versión B — divide y vencerás de Hirschberg}

El algoritmo de Hirschberg~\cite{hirschberg1975} obtiene la LCS en
tiempo $O(mn)$ y espacio $O(\min\{m,n\})$, evitando la tabla
bidimensional. La idea es:

\begin{enumerate}
    \item Partir $X$ en dos mitades $X_1$ y $X_2$ de tamaños similares.
    \item Calcular, mediante programación dinámica unidimensional, los
          vectores de longitudes $L_{1}(j)=L(|X_1|,j)$ y
          $L_{2}(j)=L(|X_2|,n-j)$.
    \item Elegir el punto de corte $k$ que maximiza $L_1(k)+L_2(k)$.
    \item Resolver recursivamente $(X_1, Y[0..k))$ y $(X_2, Y[k..n))$ y
          concatenar las dos subsecuencias obtenidas. El caso base
          ($|X|=1$) devuelve ese símbolo si aparece en $Y$.
\end{enumerate}

La recursión alcanza profundidad $O(\log m)$. En el nivel $k$ los
subproblemas suman como mucho $mn/2^{k}$ comparaciones, de modo que el
total está acotado por $mn(1+\tfrac12+\tfrac14+\dots)\le 2mn$; de ahí el
coste temporal $O(mn)$, con una constante aproximadamente doble que la
de la DP tabular.
La memoria auxiliar está dominada por los vectores de tamaño $n$, por lo
que es $O(n)$ (o $O(\min\{m,n\})$ si se traspone la cadena más corta).

\subsection{Versión C — programación dinámica tabular}

Se rellena una matriz $\mathit{dp}$ de tamaño $(m+1)\times(n+1)$
siguiendo el orden natural fila a fila. El valor final
$\mathit{dp}[m][n]$ es la longitud de la LCS. El coste temporal es
$O(mn)$ y el espacial es $O(mn)$ por la matriz completa. La
subsecuencia se reconstruye recorriendo la tabla desde $(m,n)$ hacia
$(0,0)$: en diagonal cuando los símbolos coinciden y, en otro caso, hacia
la celda vecina con mayor valor. Este recorrido cuesta $O(m+n)$. Si solo
interesara la \emph{longitud} bastaría con dos filas ($O(n)$), pero la
reconstrucción necesita la tabla completa.

\subsection{Análisis de complejidad}

La Tabla~\ref{tab:complejidad} resume las complejidades teóricas de las
tres versiones. $T(n)$ denota el caso en que $m=n$; para el caso general
basta sustituir $n$ por $\max\{m,n\}$ en la cota superior.

\begin{table}[h]
\centering
\begin{tabular}{@{}llll@{}}
\toprule
\textbf{Versión} & \textbf{Tiempo} & \textbf{Memoria} & \textbf{Observación}\\
\midrule
A · Recursiva directa    & $O(2^{m+n})$                  & $O(m+n)$ pila    & Inviable a partir de $n\approx 18$\\
B · Hirschberg           & $\Theta(mn)$                  & $O(\min\{m,n\})$ & Recursión con coste lineal por nivel\\
C · DP tabular           & $\Theta(mn)$                  & $\Theta(mn)$     & Matriz completa\\
\bottomrule
\end{tabular}
\caption{Complejidad temporal y espacial de las tres versiones.}\label{tab:complejidad}
\end{table}

\subsection{Diseño experimental}

\textbf{Datos.} Se generaron cadenas aleatorias sobre el alfabeto
$\{A,C,G,T\}$ (ADN) con un generador propio, garantizando
reproducibilidad mediante semilla fija $s=42$. Se construyeron cinco
parejas con longitudes $n\in\{50,\,100,\,250,\,500,\,1000\}$, todas con
$m=n$. Para la Versión~A, dado su coste exponencial, se emplearon
longitudes adicionales $n\in\{8,10,12,14,16,18,20\}$.

\textbf{Métricas.} Para cada ejecución se registran:
\begin{itemize}
    \item tiempo de CPU medido con \texttt{time.process\_time()}
          (excluye espera de E/S);
    \item memoria pico medida con \texttt{tracemalloc} en un subproceso
          independiente, para aislar el consumo del propio algoritmo.
\end{itemize}

\textbf{Protocolo.} Cada punto experimental se repite cinco veces y se
reporta la mediana. Se fija un \emph{timeout} de 60~s por punto, que
incluye las cinco repeticiones y la ejecución adicional de medida de
memoria. El \emph{script}
\texttt{benchmark.py} (Anexo~\ref{anexo:benchmark}) genera el CSV y las
gráficas que alimentan la Sección~\ref{sec:resultados}.

\textbf{Entorno.} Intel Xeon a 2,10\,GHz, Linux x86\_64, CPython~3.13,
sin otras cargas significativas durante la ejecución.

% =====================================================================
%  3. RESULTADOS
% =====================================================================
\section{Resultados}\label{sec:resultados}

Los valores de esta sección proceden de una ejecución de
\texttt{make run} en el entorno descrito. El Anexo~\ref{anexo:datos}
recoge el CSV completo y el Anexo~\ref{anexo:repro} las instrucciones
para reproducirlos. Los tiempos absolutos variarán de un equipo a otro;
las longitudes de LCS son idénticas en cualquier ejecución.

\subsection{Versión A — crecimiento exponencial}

La Tabla~\ref{tab:A} muestra el comportamiento de la recursión directa.
El tiempo se multiplica por un factor de entre 10 y 17 con cada
incremento de dos unidades de $n$, es decir, algo más de $3^{n}$ en la
práctica. Con $n=16$ cada llamada tarda ya 1,3\,s, y a partir de $n=18$
el punto experimental supera el \emph{timeout}.

\begin{table}[h]
\centering
\begin{tabular}{@{}rrrr@{}}
\toprule
$n$ & Tiempo CPU (mediana) & Memoria pico & LCS \\
\midrule
8    & 0{,}04\,ms        & 0{,}6\,KB & 5 \\
10   & 0{,}60\,ms        & 0{,}8\,KB & 6 \\
12   & 10\,ms            & 0{,}9\,KB & 5 \\
14   & 126\,ms           & 1{,}0\,KB & 7 \\
16   & 1{,}30\,s         & 1{,}0\,KB & 8 \\
18   & \textit{timeout}  & ---       & --- \\
20   & \textit{timeout}  & ---       & --- \\
\bottomrule
\end{tabular}
\caption{Versión A sobre cadenas de ADN de pequeña longitud.}\label{tab:A}
\end{table}

\subsection{Versiones B y C — comparación a gran escala}

La Tabla~\ref{tab:BC} contrasta Hirschberg y la DP tabular en las cinco
longitudes exigidas. Ambas resuelven el caso de $1000$ símbolos en
menos de 0,2\,s, pero difieren sustancialmente en memoria.

\begin{table}[h]
\centering
\begin{tabular}{@{}rrrrrr@{}}
\toprule
& \multicolumn{2}{c}{\textbf{B · Hirschberg}}
& \multicolumn{2}{c}{\textbf{C · DP tabular}} & \\
\cmidrule(lr){2-3}\cmidrule(lr){4-5}
$n$ & tiempo (ms) & memoria & tiempo (ms) & memoria & LCS\\
\midrule
50   & 0{,}50  & 3{,}2\,KB & 0{,}22 & 21\,KB    & 28 \\
100  & 1{,}72  & 5{,}6\,KB & 0{,}85 & 83\,KB    & 60 \\
250  & 8{,}27  & 10\,KB    & 4{,}08 & 507\,KB   & 162 \\
500  & 32{,}4  & 19\,KB    & 21{,}3 & 2{,}1\,MB & 322 \\
1000 & 159{,}4 & 66\,KB    & 80{,}8 & 11{,}5\,MB & 644 \\
\bottomrule
\end{tabular}
\caption{Comparación de las versiones B y C sobre cinco parejas de
cadenas.}\label{tab:BC}
\end{table}

La Figura~\ref{fig:tiempo-memoria} representa gráficamente ambos
indicadores en escala logarítmica, junto con la Versión~A en la gráfica
de tiempo. El tiempo de B y C crece de forma cuadrática con $n$: entre
$n=50$ y $n=1000$ la pendiente log-log es 1,9 para Hirschberg y 2,0 para
la DP tabular. La memoria sigue pendientes muy distintas: 2,1 para la DP
tabular y 1,0 para Hirschberg, coherente con $\Theta(mn)$ y
$O(\min\{m,n\})$ respectivamente.

\begin{figure}[h]
\centering
\begin{tikzpicture}
\begin{axis}[
  width=0.48\linewidth, height=5.5cm,
  xlabel={$n$}, ylabel={tiempo (ms)},
  xmode=log, ymode=log,
  legend style={at={(0.97,0.03)},anchor=south east},
  title={Tiempo de CPU},
  grid=both, grid style={gray!20},
]
\addplot[mark=triangle*,green!50!black] coordinates {(8,0.042)(10,0.602)(12,10.20)(14,125.7)(16,1298.9)};
\addplot[mark=*,blue] coordinates {(50,0.496)(100,1.718)(250,8.272)(500,32.37)(1000,159.4)};
\addplot[mark=square*,red] coordinates {(50,0.218)(100,0.848)(250,4.080)(500,21.29)(1000,80.84)};
\legend{A, B, C}
\end{axis}
\end{tikzpicture}\hfill
\begin{tikzpicture}
\begin{axis}[
  width=0.48\linewidth, height=5.5cm,
  xlabel={$n$}, ylabel={memoria pico (KB)},
  xmode=log, ymode=log,
  legend style={at={(0.03,0.97)},anchor=north west},
  title={Memoria máxima},
  grid=both, grid style={gray!20},
]
\addplot[mark=*,blue] coordinates {(50,3.15)(100,5.55)(250,10.26)(500,18.98)(1000,66.46)};
\addplot[mark=square*,red] coordinates {(50,21.31)(100,82.97)(250,506.7)(500,2140.5)(1000,11736.3)};
\legend{B, C}
\end{axis}
\end{tikzpicture}
\caption{Tiempo de CPU y memoria máxima en función de $n$
(escala log-log).}\label{fig:tiempo-memoria}
\end{figure}

% =====================================================================
%  4. DISCUSIÓN
% =====================================================================
\section{Discusión}\label{sec:discusion}

\paragraph{Recursión directa: coste teórico confirmado.}
Los datos de la Tabla~\ref{tab:A} verifican el comportamiento
exponencial predicho por el análisis. El tiempo se multiplica por unas
diez veces cada dos unidades de $n$, lo que refleja la estructura de
árbol binario de la recurrencia. El crecimiento observado queda por
debajo de la cota $2^{m+n}=4^{n}$ porque cada coincidencia de símbolos
genera una sola llamada en lugar de dos; con alfabeto de cuatro letras,
aproximadamente una de cada cuatro comparaciones es una coincidencia.
Para $n\le 14$ la versión responde en menos de 0,2\,s, con $n=18$ ya no
completa las seis ejecuciones en 60\,s, y para $n=1000$ sería
inabordable.

La memoria pico medida apenas pasa de 0,6 a 1,0\,KB, y corresponde a las
cadenas parciales que se construyen en cada nivel. \texttt{tracemalloc}
registra los objetos creados en el \emph{heap}, pero no los marcos de la
pila de recursión, que CPython~3.11+ gestiona en una zona propia. El
consumo real crece linealmente con la profundidad de la pila, que nunca
supera $m+n$ llamadas anidadas, de acuerdo con la cota $O(m+n)$.

\paragraph{Hirschberg frente a DP tabular.}
Ambas versiones comparten la cota temporal $\Theta(mn)$, pero
Hirschberg es entre 1,5 y 2,3 veces más lenta según el tamaño. La causa
es la que anticipa el análisis: el algoritmo recorre del orden de $2mn$
celdas frente a las $mn$ de la DP tabular, y además paga las llamadas
recursivas, las copias de subcadenas y la inversión de cadenas en cada
nivel. La diferencia se vuelve decisiva en memoria: para $n=1000$, la
tabla completa ocupa 11,5\,MB, mientras que Hirschberg necesita 66\,KB,
unas 180 veces menos. Las dos versiones obtienen siempre la misma
longitud, aunque no necesariamente la misma subsecuencia: con $n=1000$
devuelven cadenas distintas de 644 símbolos, ambas óptimas, porque el
problema admite varias soluciones y cada algoritmo deshace los empates
de forma diferente.

El salto de memoria de Hirschberg entre $n=500$ (18\,KB) y $n=1000$
(66\,KB) es mayor que el factor 2 esperado. Se explica por un detalle de
CPython: los enteros entre $-5$ y $256$ están precalculados y no ocupan
memoria nueva, mientras que los mayores se crean como objetos de 28
bytes. Con $n=1000$ los valores de los vectores de longitudes superan 256
y cada celda pasa a reservar un objeto propio. La pendiente asintótica
sigue siendo lineal. El mismo efecto contribuye a que la memoria de la DP
tabular crezca algo más deprisa que $n^{2}$ en el último tramo.

\paragraph{¿Cuándo elegir cada versión?}
La elección depende del recurso crítico:
\begin{itemize}
    \item Si solo se necesita la \emph{longitud}, basta una DP con dos
          filas, que es la más simple y rápida.
    \item Si se necesita la subsecuencia y $mn$ cabe en memoria, la DP
          tabular con reconstrucción es la opción más directa.
    \item Si se necesita reconstruir la subsecuencia y las cadenas son
          grandes (genómica), Hirschberg es la opción profesional:
          mismo coste temporal, memoria lineal.
    \item La versión recursiva directa tiene valor didáctico, pero no de
          producción.
\end{itemize}

\paragraph{Limitaciones del experimento.}
Las mediciones se tomaron sobre un único equipo, con cadenas aleatorias
de alfabeto uniforme. Los tiempos absolutos dependen fuertemente del
lenguaje y de la implementación: en Python el intérprete domina,
mientras que en C/C++ las diferencias entre B y C serían menores en
tiempo. La estructura de los datos (por ejemplo, alfabetos muy pequeños
o cadenas con muchos símbolos repetidos) también altera el factor
constante, aunque no la complejidad asintótica.

% =====================================================================
%  5. CONCLUSIONES
% =====================================================================
\section{Conclusiones}\label{sec:conclusiones}

Se han diseñado, analizado y comparado empíricamente tres algoritmos
para el problema de la LCS. Los experimentos confirman:

\begin{enumerate}
    \item que la \textbf{recursión directa} incurre en un coste
          exponencial que la invalida a partir de $n\approx 18$ en el
          entorno de pruebas, tal y como predice el análisis;
    \item que la \textbf{estrategia de Hirschberg} alcanza el mismo
          orden temporal que la DP tabular con un consumo de memoria dos
          órdenes de magnitud menor (66\,KB frente a 11,5\,MB con
          $n=1000$), a cambio de duplicar aproximadamente el tiempo, y es
          la opción preferente cuando el espacio es un recurso crítico;
    \item que la \textbf{DP tabular} es la solución más directa y la de
          menor coste constante, y resulta adecuada cuando la tabla cabe
          en memoria.
\end{enumerate}

Como trabajo futuro se propone extender el estudio a cadenas con
estructura real (por ejemplo, secuencias genómicas), añadir una versión
memorizada (\emph{top-down}) de la recursión, y repetir los
experimentos en un lenguaje compilado para cuantificar el peso del
intérprete.

% =====================================================================
%  REFERENCIAS
% =====================================================================
\begin{thebibliography}{9}
\bibitem{hirschberg1975} D.~S. Hirschberg, \emph{A linear space
algorithm for computing maximal common subsequences}, Communications of
the ACM, 18(6):341--343, 1975.
\bibitem{clrs} T.~H. Cormen, C.~E. Leiserson, R.~L. Rivest y
C.~Stein, \emph{Introduction to Algorithms}, 3.\textsuperscript{a} ed.,
MIT Press, 2009, cap.~15.
\bibitem{sankoff} D. Sankoff y J.~B. Kruskal, \emph{Time Warps, String
Edits, and Macromolecules: The Theory and Practice of Sequence
Comparison}, Addison-Wesley, 1983.
\end{thebibliography}

% =====================================================================
%  ANEXOS
% =====================================================================
\clearpage
\appendix
\pagenumbering{Alph}

% ---------------------------------------------------------------------
\section{Código fuente completo}\label{anexo:codigo}

Este anexo contiene la totalidad del código entregado. Los ficheros
están organizados de modo que, copiados en una carpeta local con la
estructura indicada en el Anexo~\ref{anexo:repro}, el proyecto se
ejecuta con un único comando (\texttt{make run}). Todo el contenido está
incrustado en este documento y también está disponible en el repositorio
\url{https://github.com/carpediemservus/lcs-actividad1}.

\subsection{Estructura del proyecto}

\begin{verbatim}
.
|-- Makefile
|-- README.md
|-- src/
|   |-- lcs.py          # implementaciones A, B, C
|   |-- benchmark.py    # experimentos + CSV + gráficas
|   |-- test_lcs.py     # pruebas de correctitud
|   `-- analyze.py      # tablas LaTeX a partir del CSV
`-- results/            # generado por `make run`
    |-- resultados.csv
    |-- tiempo.png
    |-- memoria.png
    `-- datos/          # parejas X, Y y LCS de cada n
\end{verbatim}

\lstlistoflistings

\subsection{\texttt{src/lcs.py} — implementaciones A, B y C}\label{anexo:lcs}

\begin{lstlisting}[style=py, caption={\texttt{src/lcs.py}}]
"""Tres implementaciones de la Subsecuencia Comun Mas Larga (LCS).

Version A: recursion directa, sin memorizacion.
Version B: divide y venceras de Hirschberg (memoria lineal).
Version C: programacion dinamica tabular con reconstruccion.

Las tres devuelven la propia subsecuencia (una cadena); su longitud se
obtiene con len().
"""
import sys

sys.setrecursionlimit(10000)


# ---------- Version A: recursiva directa ----------
def lcs_recursiva(a: str, b: str) -> str:
    def rec(i: int, j: int) -> str:
        if i == 0 or j == 0:
            return ""
        if a[i - 1] == b[j - 1]:
            return rec(i - 1, j - 1) + a[i - 1]
        x = rec(i - 1, j)
        y = rec(i, j - 1)
        return x if len(x) >= len(y) else y

    return rec(len(a), len(b))


# ---------- Version B: Hirschberg ----------
def _last_row(a: str, b: str) -> list:
    """Ultima fila de la DP (longitudes LCS(a, b[:j])) en O(|b|) memoria."""
    n = len(b)
    prev = [0] * (n + 1)
    for ca in a:
        curr = [0] * (n + 1)
        for j in range(1, n + 1):
            if ca == b[j - 1]:
                curr[j] = prev[j - 1] + 1
            else:
                curr[j] = prev[j] if prev[j] >= curr[j - 1] else curr[j - 1]
        prev = curr
    return prev


def lcs_hirschberg(a: str, b: str) -> str:
    if not a or not b:
        return ""
    if len(a) == 1:
        return a if a in b else ""

    mid = len(a) // 2
    l1 = _last_row(a[:mid], b)
    l2 = _last_row(a[mid:][::-1], b[::-1])[::-1]

    # punto de corte k que maximiza L1(k) + L2(k)
    k, best = 0, -1
    for j in range(len(b) + 1):
        s = l1[j] + l2[j]
        if s > best:
            best, k = s, j

    return lcs_hirschberg(a[:mid], b[:k]) + lcs_hirschberg(a[mid:], b[k:])


# ---------- Version C: DP tabular ----------
def lcs_dp(a: str, b: str) -> str:
    m, n = len(a), len(b)
    dp = [[0] * (n + 1) for _ in range(m + 1)]
    for i in range(1, m + 1):
        ai = a[i - 1]
        row, prev = dp[i], dp[i - 1]
        for j in range(1, n + 1):
            if ai == b[j - 1]:
                row[j] = prev[j - 1] + 1
            else:
                row[j] = prev[j] if prev[j] >= row[j - 1] else row[j - 1]

    # reconstruccion: recorrido inverso desde dp[m][n]
    out = []
    i, j = m, n
    while i > 0 and j > 0:
        if a[i - 1] == b[j - 1]:
            out.append(a[i - 1])
            i -= 1
            j -= 1
        elif dp[i - 1][j] >= dp[i][j - 1]:
            i -= 1
        else:
            j -= 1
    return "".join(reversed(out))
\end{lstlisting}

\subsection{\texttt{src/test\_lcs.py} — pruebas de correctitud}\label{anexo:test}

\begin{lstlisting}[style=py, caption={\texttt{src/test\_lcs.py}}]
"""Pruebas de correctitud de las tres versiones.

Para cada caso se comprueba que la cadena devuelta:
  1. es subsecuencia de ambas entradas, y
  2. tiene la longitud optima, calculada con una DP de referencia.
"""
import random

from lcs import lcs_recursiva, lcs_hirschberg, lcs_dp


def _reference_len(a: str, b: str) -> int:
    m, n = len(a), len(b)
    dp = [[0] * (n + 1) for _ in range(m + 1)]
    for i in range(1, m + 1):
        for j in range(1, n + 1):
            if a[i - 1] == b[j - 1]:
                dp[i][j] = dp[i - 1][j - 1] + 1
            else:
                dp[i][j] = max(dp[i - 1][j], dp[i][j - 1])
    return dp[m][n]


def es_subsecuencia(s: str, t: str) -> bool:
    it = iter(t)
    return all(c in it for c in s)


def run(trials: int = 300) -> None:
    rng = random.Random(0)
    cases = [("", ""), ("A", ""), ("", "A"), ("A", "A"), ("AB", "BA"),
             ("ACGT", "ACGT"), ("AAAA", "CCCC"), ("ABCBDAB", "BDCABA")]
    for _ in range(trials):
        a = "".join(rng.choice("ACGT") for _ in range(rng.randint(0, 12)))
        b = "".join(rng.choice("ACGT") for _ in range(rng.randint(0, 12)))
        cases.append((a, b))
    # algunos casos mayores solo para B y C
    big = []
    for n in (100, 300):
        a = "".join(rng.choice("ACGT") for _ in range(n))
        b = "".join(rng.choice("ACGT") for _ in range(n))
        big.append((a, b))

    versiones = [("A", lcs_recursiva), ("B", lcs_hirschberg), ("C", lcs_dp)]
    total = 0
    for a, b in cases + big:
        expected = _reference_len(a, b)
        for name, fn in versiones:
            if name == "A" and (a, b) in big:
                continue
            got = fn(a, b)
            assert len(got) == expected, (
                f"Version {name}: {a!r} {b!r} -> longitud {len(got)} != {expected}"
            )
            assert es_subsecuencia(got, a) and es_subsecuencia(got, b), (
                f"Version {name}: {got!r} no es subsecuencia comun de {a!r} y {b!r}"
            )
            total += 1
    print(f"OK: {len(cases)} casos en las 3 versiones y {len(big)} casos "
          f"grandes en B y C ({total} comprobaciones).")


if __name__ == "__main__":
    run()
\end{lstlisting}

\subsection{\texttt{src/benchmark.py} — protocolo experimental}\label{anexo:benchmark}

\begin{lstlisting}[style=py, caption={\texttt{src/benchmark.py}}]
"""Benchmarks de las tres versiones de LCS.

Para cada (version, n) lanza un subproceso que:
  - genera las cadenas con semilla fija,
  - calcula la LCS (la subsecuencia) y registra su longitud,
  - mide el tiempo de CPU (mediana de REPEATS),
  - mide la memoria pico de Python con tracemalloc.
El proceso padre aplica timeout, escribe un CSV, dos graficas (log-log) y
guarda las parejas de cadenas y las LCS obtenidas en results/datos/.

Uso:
    python3 src/benchmark.py --csv results/resultados.csv
"""
import argparse
import csv
import json
import random
import statistics
import subprocess
import sys
import time
import tracemalloc
from pathlib import Path

HERE = Path(__file__).resolve().parent
sys.path.insert(0, str(HERE))

ALPHABET = "ACGT"
SEED = 42
REPEATS = 5
TIMEOUT = 60  # segundos

CONFIGS = {
    "A": [8, 10, 12, 14, 16, 18, 20],
    "B": [50, 100, 250, 500, 1000],
    "C": [50, 100, 250, 500, 1000],
}


# ----------------------- worker -----------------------
def _gen(n: int, seed: int) -> tuple:
    rng = random.Random(seed)
    return (
        "".join(rng.choice(ALPHABET) for _ in range(n)),
        "".join(rng.choice(ALPHABET) for _ in range(n)),
    )


def worker(version: str, n: int, seed: int) -> None:
    from lcs import lcs_recursiva, lcs_hirschberg, lcs_dp
    fn = {"A": lcs_recursiva, "B": lcs_hirschberg, "C": lcs_dp}[version]

    a, b = _gen(n, seed)

    # tiempo (mediana de REPEATS)
    times, result = [], None
    for _ in range(REPEATS):
        t0 = time.process_time()
        result = fn(a, b)
        times.append(time.process_time() - t0)

    # memoria pico (asignaciones Python durante la llamada)
    tracemalloc.start()
    fn(a, b)
    peak = tracemalloc.get_traced_memory()[1]
    tracemalloc.stop()

    print(json.dumps({
        "version": version,
        "n": n,
        "time_s": statistics.median(times),
        "mem_kb": peak / 1024.0,
        "lcs": len(result),
        "sub": result,
    }))


# ----------------------- driver -----------------------
def run_one(version: str, n: int) -> dict:
    cmd = [sys.executable, str(Path(__file__).resolve()),
           "--worker", version, str(n), str(SEED)]
    try:
        out = subprocess.run(cmd, capture_output=True, text=True,
                             timeout=TIMEOUT, check=True)
        return json.loads(out.stdout)
    except subprocess.TimeoutExpired:
        return {"version": version, "n": n, "time_s": None,
                "mem_kb": None, "lcs": None, "status": "timeout"}


def write_csv(rows, path: Path) -> None:
    path.parent.mkdir(parents=True, exist_ok=True)
    with path.open("w", newline="") as f:
        w = csv.writer(f)
        w.writerow(["version", "n", "time_ms", "mem_kb", "lcs", "status"])
        for r in rows:
            t = "" if r.get("time_s") is None else f"{r['time_s']*1000:.3f}"
            m = "" if r.get("mem_kb") is None else f"{r['mem_kb']:.2f}"
            w.writerow([r["version"], r["n"], t, m,
                        r.get("lcs", ""), r.get("status", "ok")])


def write_data(rows, outdir: Path) -> None:
    """Guarda cada pareja (X, Y) y la LCS obtenida por la version C."""
    d = outdir / "datos"
    d.mkdir(parents=True, exist_ok=True)
    for r in rows:
        if r["version"] != "C" or r.get("sub") is None:
            continue
        x, y = _gen(r["n"], SEED)
        (d / f"n{r['n']}.txt").write_text(
            f"X={x}\nY={y}\nLCS={r['sub']}\n", encoding="utf-8")
    print(f"[benchmark] datos -> {d}")


def write_plots(rows, outdir: Path) -> None:
    try:
        import matplotlib
        matplotlib.use("Agg")
        import matplotlib.pyplot as plt
    except ImportError:
        print("[benchmark] matplotlib no instalado: se omiten las graficas.")
        return

    outdir.mkdir(parents=True, exist_ok=True)
    colors = {"A": "tab:green", "B": "tab:blue", "C": "tab:red"}
    marks = {"A": "^", "B": "o", "C": "s"}

    for field, fname, ylabel, title in [
        ("time_s", "tiempo.png", "tiempo CPU (s)", "Tiempo de CPU"),
        ("mem_kb", "memoria.png", "memoria pico (KB)", "Memoria maxima"),
    ]:
        fig, ax = plt.subplots(figsize=(5.5, 4))
        for v in ("A", "B", "C"):
            pts = [(r["n"], r[field]) for r in rows
                   if r["version"] == v and r.get(field) is not None]
            if not pts:
                continue
            xs, ys = zip(*sorted(pts))
            ax.plot(xs, ys, marker=marks[v], color=colors[v],
                    label=f"Version {v}")
        ax.set_xscale("log")
        ax.set_yscale("log")
        ax.set_xlabel("n")
        ax.set_ylabel(ylabel)
        ax.set_title(title)
        ax.grid(True, which="both", ls=":", alpha=0.4)
        ax.legend()
        fig.tight_layout()
        fig.savefig(outdir / fname, dpi=150)
        plt.close(fig)
        print(f"[benchmark] grafica -> {outdir / fname}")


# ----------------------- main -----------------------
def main() -> None:
    ap = argparse.ArgumentParser()
    ap.add_argument("--worker", nargs=3, metavar=("VERSION", "N", "SEED"))
    ap.add_argument("--csv", default="results/resultados.csv")
    args = ap.parse_args()

    if args.worker:
        v, n, s = args.worker
        worker(v, int(n), int(s))
        return

    rows = []
    for version, sizes in CONFIGS.items():
        for n in sizes:
            print(f"[benchmark] version {version}, n={n} ...", flush=True)
            r = run_one(version, n)
            rows.append(r)
            if r.get("status") == "timeout":
                print("   -> timeout")
            else:
                print(f"   -> t={r['time_s']*1000:8.2f} ms   "
                      f"mem={r['mem_kb']:8.1f} KB   LCS={r['lcs']}")

    csv_path = Path(args.csv)
    write_csv(rows, csv_path)
    print(f"[benchmark] CSV -> {csv_path}")
    write_data(rows, csv_path.parent)
    write_plots(rows, csv_path.parent)


if __name__ == "__main__":
    main()
\end{lstlisting}

\subsection{\texttt{src/analyze.py} — generación de tablas LaTeX}\label{anexo:analyze}

\begin{lstlisting}[style=py, caption={\texttt{src/analyze.py}}]
"""Convierte results/resultados.csv en las Tablas 2 y 3 del informe (LaTeX).

Uso:
    python3 src/analyze.py
    python3 src/analyze.py --csv otro.csv
    python3 src/analyze.py --out tablas.tex
"""
import argparse
import csv
from pathlib import Path


def fmt_time(ms):
    """Tiempo en ms, con 2 decimales y separador decimal LaTeX (coma)."""
    if ms is None or ms == "":
        return "---"
    v = float(ms)
    if v >= 1000:
        s = f"{v/1000:.2f}".rstrip("0").rstrip(".")
        s = s.replace(".", "{,}")
        return f"{s}\\,s"
    s = f"{v:.2f}".rstrip("0").rstrip(".")
    return s.replace(".", "{,}")


def fmt_mem(kb):
    """Memoria: KB si <1024, MB en otro caso (con 2 decimales)."""
    if kb is None or kb == "":
        return "---"
    v = float(kb)
    if v >= 1024:
        s = f"{v/1024:.1f}".rstrip("0").rstrip(".")
        return s.replace(".", "{,}") + "\\,MB"
    s = f"{v:.1f}" if v < 10 else f"{v:.0f}"
    return s.replace(".", "{,}") + "\\,KB"


def fmt_lcs(v):
    return "---" if v in (None, "") else str(v)


def load(csv_path):
    rows = []
    with open(csv_path, newline="") as f:
        for r in csv.DictReader(f):
            rows.append({
                "version": r["version"],
                "n": int(r["n"]),
                "time_ms": None if r["time_ms"] == "" else float(r["time_ms"]),
                "mem_kb":  None if r["mem_kb"] == ""  else float(r["mem_kb"]),
                "lcs":     None if r["lcs"] == ""     else r["lcs"],
                "status":  r.get("status", "ok"),
            })
    return rows


def by_version(rows, v):
    return sorted([r for r in rows if r["version"] == v], key=lambda r: r["n"])


def table_A(rows):
    out = [
        r"\begin{table}[h]", r"\centering",
        r"\begin{tabular}{@{}rrrr@{}}", r"\toprule",
        r"$n$ & Tiempo CPU (mediana) & Memoria pico & LCS \\",
        r"\midrule",
    ]
    for r in by_version(rows, "A"):
        t = "\\textit{timeout}" if r["status"] == "timeout" else fmt_time(r["time_ms"])
        m = fmt_mem(r["mem_kb"])
        l = fmt_lcs(r["lcs"])
        out.append(f"{r['n']:<4} & {t:<20} & {m:<10} & {l} \\\\")
    out += [r"\bottomrule", r"\end{tabular}",
            r"\caption{Version A.}\label{tab:A}", r"\end{table}"]
    return "\n".join(out)


def table_BC(rows):
    b = {r["n"]: r for r in by_version(rows, "B")}
    c = {r["n"]: r for r in by_version(rows, "C")}
    ns = sorted(set(b) | set(c))
    out = [
        r"\begin{table}[h]", r"\centering",
        r"\begin{tabular}{@{}rrrrrr@{}}", r"\toprule",
        r"& \multicolumn{2}{c}{\textbf{B - Hirschberg}} "
        r"& \multicolumn{2}{c}{\textbf{C - DP tabular}} & \\",
        r"\cmidrule(lr){2-3}\cmidrule(lr){4-5}",
        r"$n$ & tiempo (ms) & memoria & tiempo (ms) & memoria & LCS\\",
        r"\midrule",
    ]
    for n in ns:
        rb, rc = b.get(n, {}), c.get(n, {})
        out.append(
            f"{n:<5} & {fmt_time(rb.get('time_ms')):<10} & "
            f"{fmt_mem(rb.get('mem_kb')):<9} & "
            f"{fmt_time(rc.get('time_ms')):<10} & "
            f"{fmt_mem(rc.get('mem_kb')):<9} & "
            f"{fmt_lcs(rb.get('lcs') or rc.get('lcs'))} \\\\"
        )
    out += [r"\bottomrule", r"\end{tabular}",
            r"\caption{Versiones B y C.}\label{tab:BC}", r"\end{table}"]
    return "\n".join(out)


def main():
    ap = argparse.ArgumentParser()
    ap.add_argument("--csv", default="results/resultados.csv")
    ap.add_argument("--out", default=None)
    args = ap.parse_args()

    rows = load(Path(args.csv))
    text = ("% === Tabla 2 ===\n" + table_A(rows)
            + "\n\n% === Tabla 3 ===\n" + table_BC(rows) + "\n")
    if args.out:
        Path(args.out).write_text(text, encoding="utf-8")
        print(f"Tablas escritas en {args.out}")
    else:
        print(text)


if __name__ == "__main__":
    main()
\end{lstlisting}

\subsection{\texttt{Makefile}}\label{anexo:makefile}

\begin{lstlisting}[style=mk, caption={\texttt{Makefile}}]
PYTHON ?= python3
CSV    := results/resultados.csv

.PHONY: all run test analyze clean help

all: run

run:
	@mkdir -p results
	$(PYTHON) src/benchmark.py --csv $(CSV)

test:
	$(PYTHON) src/test_lcs.py

analyze: run
	$(PYTHON) src/analyze.py --csv $(CSV) --out results/tablas.tex
	@echo "Tablas generadas en results/tablas.tex"

clean:
	rm -rf results
	find . -type d -name __pycache__ -exec rm -rf {} + 2>/dev/null || true

help:
	@echo "make run     - ejecuta benchmarks (CSV + graficas)"
	@echo "make test    - pruebas de correctitud"
	@echo "make analyze - vuelca Tablas 2 y 3 en LaTeX"
	@echo "make clean   - limpia artefactos"
\end{lstlisting}

% ---------------------------------------------------------------------
\section{Datos y resultados brutos}\label{anexo:datos}

\subsection{Generación de los datos}

Todas las cadenas empleadas en el informe se generan de forma
determinista con el siguiente procedimiento (función \texttt{\_gen} del
Anexo~\ref{anexo:benchmark}):

\begin{itemize}
    \item Alfabeto: $\{A,C,G,T\}$.
    \item Semilla: $s=42$.
    \item Para cada $n$, se construyen dos cadenas independientes
          $X,Y\in\{A,C,G,T\}^n$ usando un generador pseudoaleatorio
          \texttt{random.Random(42)}.
\end{itemize}

Al fijar la semilla, cualquier ejecución de \texttt{make run} reproduce
exactamente las mismas cadenas y, por tanto, los mismos valores de LCS.
Las cinco parejas completas, junto con la LCS obtenida por la
Versión~C, se guardan en \texttt{results/datos/n<n>.txt} y están
publicadas en el repositorio. Como muestra, la pareja de $n=50$ es:

\begin{lstlisting}[style=txt, caption={\texttt{results/datos/n50.txt}}]
X=AAGCCCAATAAACCACTCTGACTGGCCGAATAGGGATATAGGCAACGACA
Y=TGTGCGGCGACCCTTGCGACAGTGACGCTTTCGCCGTTGCCTAAACCTAT
LCS=GCCACCCTCGACTGGCCGGTTGCAACCA
\end{lstlisting}
La Tabla~\ref{tab:lcs-valores} recoge la longitud de la LCS obtenida
para cada par. Las versiones B y C coinciden en todos los casos, y el
cociente LCS$/n$ tiende a 0,64--0,65, el valor esperado para cadenas
aleatorias sobre un alfabeto de cuatro símbolos.

\begin{table}[h]
\centering
\begin{tabular}{@{}lrrrrr@{}}
\toprule
$n$ & 50 & 100 & 250 & 500 & 1000 \\
\midrule
LCS & 28 & 60 & 162 & 322 & 644 \\
\bottomrule
\end{tabular}
\caption{Longitud de la LCS para cada pareja de cadenas (semilla
$s=42$).}\label{tab:lcs-valores}
\end{table}

\subsection{Volcado del CSV de resultados}

El Anexo~\ref{anexo:benchmark} genera el siguiente fichero
\texttt{results/resultados.csv}. Se incluye aquí como evidencia de los
tiempos y memoria reportados en la Sección~\ref{sec:resultados}.

\begin{lstlisting}[style=txt, caption={\texttt{results/resultados.csv}}]
version,n,time_ms,mem_kb,lcs,status
A,8,0.042,0.64,5,ok
A,10,0.602,0.82,6,ok
A,12,10.203,0.90,5,ok
A,14,125.731,0.99,7,ok
A,16,1298.861,1.04,8,ok
A,18,,,,timeout
A,20,,,,timeout
B,50,0.496,3.15,28,ok
B,100,1.718,5.55,60,ok
B,250,8.272,10.26,162,ok
B,500,32.370,18.98,322,ok
B,1000,159.369,66.46,644,ok
C,50,0.218,21.31,28,ok
C,100,0.848,82.97,60,ok
C,250,4.080,506.66,162,ok
C,500,21.289,2140.45,322,ok
C,1000,80.841,11736.25,644,ok
\end{lstlisting}

\subsection{Pruebas de correctitud}

La ejecución de \texttt{make test} (Anexo~\ref{anexo:test}) comprueba,
para cada caso, que la cadena devuelta es subsecuencia de las dos
entradas y que su longitud coincide con la óptima calculada por una DP
de referencia. Se usan 308 casos con las tres versiones (ocho casos
borde y 300 pares aleatorios de longitud entre 0 y 12) y dos pares de
100 y 300 símbolos solo con B y C, en total 928 comprobaciones. Todas
finalizan con éxito:

\begin{lstlisting}[style=txt, caption={Salida de \texttt{make test}}]
OK: 308 casos en las 3 versiones y 2 casos grandes en B y C (928 comprobaciones).
\end{lstlisting}

% ---------------------------------------------------------------------
\section{Instrucciones de reproducción local}\label{anexo:repro}

Todo el código y los datos necesarios están contenidos en este
documento y en el repositorio público
\url{https://github.com/carpediemservus/lcs-actividad1}. Para reproducir
los resultados en un equipo local:

\begin{enumerate}
    \item Clone el repositorio:
          \texttt{git clone https://github.com/carpediemservus/lcs-actividad1}.
          Como alternativa, cree una carpeta de trabajo y copie en ella los
          ficheros del Anexo~\ref{anexo:codigo} respetando la estructura
          indicada (carpeta \texttt{src/} y \texttt{Makefile} en la raíz).
    \item Verifique que dispone de Python~3.9 o superior:
          \texttt{python3 --version}.
    \item (Opcional) Instale \texttt{matplotlib} si desea regenerar las
          gráficas: \texttt{pip install matplotlib}.
    \item Ejecute:
\end{enumerate}

\begin{verbatim}
make test      # valida las 3 versiones contra una DP de referencia
make run       # ejecuta los benchmarks, genera CSV y graficas
make analyze   # vuelca las Tablas 2 y 3 en LaTeX
\end{verbatim}

Los artefactos quedan en \texttt{results/}:

\begin{itemize}
    \item \texttt{resultados.csv} — datos brutos de todas las mediciones;
    \item \texttt{tiempo.png}, \texttt{memoria.png} — gráficas log-log;
    \item \texttt{tablas.tex} — Tablas 2 y 3 listas para pegar en el
          informe.
\end{itemize}

El tiempo total de ejecución de \texttt{make run} es de
aproximadamente 2--3 minutos, de los que 120~s corresponden a los dos
\emph{timeouts} de la Versión~A ($n=18$ y $n=20$). El compilado del informe no requiere ningún paso
adicional más allá de \texttt{pdflatex informe.tex} dos veces.

\end{document}